% --- Archon physics formalization source begin ---
% archon:physics
% archon:covers PhyXMiniProblems/problem_phyx_mini_0589.lean
% archon:source-report reports/phyx_mini/problem_phyx_mini_0589.source.json

\chapter{Physics problem phyx\_mini\_0589}
\label{ch:PhyXMiniProblems_problem_phyx_mini_0589}

\paragraph{Problem source.}
\#\# Physical scenario

In figure, two cruisers fly toward a space station. Relative to the station, cruiser \textbackslash{}( A \textbackslash{}) has speed \textbackslash{}( 0.800c \textbackslash{}).

\#\# Question

Relative to the station, what speed is required of cruiser \textbackslash{}( B \textbackslash{}) such that its pilot sees \textbackslash{}( A \textbackslash{}) and the station approach \textbackslash{}( B \textbackslash{}) at the same speed?

\#\# Answer choices

- A: A: 0.6c

- B: B: 0.4c

- C: C: 0.5c

- D: D: 0.7c

\#\# Figure caption (auxiliary; use the image as primary evidence)

The image depicts a simple physics diagram involving velocity and motion. It includes the following elements:

1. **Objects:**
   - Two rockets: One is red, and the other is blue. Both are traveling to the right.
   
2. **Arrows and Velocity:**
   - A magenta arrow labeled \textbackslash{}( \textbackslash{}vec\{v\}\_A \textbackslash{}) is above the red rocket, indicating its velocity in the direction it is moving.
   - Another magenta arrow labeled \textbackslash{}( \textbackslash{}vec\{v\}\_B \textbackslash{}) is above the blue rocket, similarly indicating its velocity.

3. **Station:**
   - There is a structure labeled "Station" on the right, depicted with two vertical blue bars and a dashed line in between.

4. **Line and Axis:**
   - A horizontal line represents an axis, labeled with "x" on the far right.

5. **Relationships:**
   - Both rockets are moving towards the right, indicated by the direction of the velocity vectors.
   - The "Station" is positioned further to the right along the horizontal axis.

The diagram conveys a scenario involving two rockets moving towards a stationary point, which might be used to illustrate concepts in physics related to relative motion or velocity.

\#\# Dataset metadata

Relativity; Physical Model Grounding Reasoning, Multi-Formula Reasoning

\paragraph{Recorded answer/context.}
C: C: 0.5c

\paragraph{Figure/image path.}
/root/proposal\_for\_physic/science-mango/phyx\_mini\_run/phyx\_data/test\_image/589.png

\paragraph{Formalization target.}
create a compiling Lean file with sorry bodies at `PhyXMiniProblems/problem\_phyx\_mini\_0589.lean`. The Lean declarations must preserve the physical quantities, dimensions or dimensional roles, figure labels, governing-law hypotheses, and final relation expressed by this problem.
Use Mathlib/Physlib names found through LeanExplore where available. If a physics API is missing, introduce faithful local abstractions rather than scalar placeholder aliases.

\begin{theorem}[Physics formalization target]
\label{thm:physics:phyx_mini_0589:target}
\lean{PhyXMiniProblems.ProblemPhyXMini0589.problem_phyx_mini_0589}
\uses{def:physics:phyx-mini-0589:phyxminiproblems-problemphyxmini0589-speedfractionoflight, def:physics:phyx-mini-0589:phyxminiproblems-problemphyxmini0589-equalapproachcruisersetup, def:physics:phyx-mini-0589:phyxminiproblems-problemphyxmini0589-matchesprimaryfigure, def:physics:phyx-mini-0589:phyxminiproblems-problemphyxmini0589-matchesproblemdata, def:physics:phyx-mini-0589:phyxminiproblems-problemphyxmini0589-hasphysicalspeedparameters, def:physics:phyx-mini-0589:phyxminiproblems-problemphyxmini0589-representsapproachdirections, def:physics:phyx-mini-0589:phyxminiproblems-problemphyxmini0589-satisfiescollineareinsteinvelocitytransformations, def:physics:phyx-mini-0589:phyxminiproblems-problemphyxmini0589-approachspeedsarevelocitymagnitudes, def:physics:phyx-mini-0589:phyxminiproblems-problemphyxmini0589-pilotseesequalapproachspeeds, def:physics:phyx-mini-0589:phyxminiproblems-problemphyxmini0589-recordeddatasetanswer, def:physics:phyx-mini-0589:phyxminiproblems-problemphyxmini0589-isuniquematchinganswerchoice}
The assigned autoformalize agent should translate this physics problem into Lean declarations in the covered file, with theorem and lemma proof bodies written as `by sorry`.
\end{theorem}
\begin{proof}
This is an autoformalization task, not a proof task. Produce statements that can later be proved by the physics prover without weakening the physical model.
\end{proof}

% --- Archon declaration topology begin ---
\section{Declaration topology}
\begin{definition}[Speed Quantity]
  \label{def:physics:phyx-mini-0589:phyxminiproblems-problemphyxmini0589-speedquantity}
  \lean{PhyXMiniProblems.ProblemPhyXMini0589.SpeedQuantity}
  A nonnegative physical speed magnitude, independent of a unit system
\end{definition}
\begin{proof}
  This is the declared model definition of Speed Quantity.
\end{proof}
\begin{definition}[Signed Velocity Quantity]
  \label{def:physics:phyx-mini-0589:phyxminiproblems-problemphyxmini0589-signedvelocityquantity}
  \lean{PhyXMiniProblems.ProblemPhyXMini0589.SignedVelocityQuantity}
  A signed one-dimensional physical velocity of dimension length per time
\end{definition}
\begin{proof}
  This is the declared model definition of Signed Velocity Quantity.
\end{proof}
\begin{definition}[speed Readout]
  \label{def:physics:phyx-mini-0589:phyxminiproblems-problemphyxmini0589-speedreadout}
  \lean{PhyXMiniProblems.ProblemPhyXMini0589.speedReadout}
  \uses{def:physics:phyx-mini-0589:phyxminiproblems-problemphyxmini0589-speedquantity}
  Read a physical speed magnitude in selected length and time units
\end{definition}
\begin{proof}
  This is the declared model definition of speed Readout.
\end{proof}
\begin{definition}[signed Velocity Readout]
  \label{def:physics:phyx-mini-0589:phyxminiproblems-problemphyxmini0589-signedvelocityreadout}
  \lean{PhyXMiniProblems.ProblemPhyXMini0589.signedVelocityReadout}
  \uses{def:physics:phyx-mini-0589:phyxminiproblems-problemphyxmini0589-signedvelocityquantity}
  Read a signed axial velocity in selected length and time units
\end{definition}
\begin{proof}
  This is the declared model definition of signed Velocity Readout.
\end{proof}
\begin{definition}[vacuum Speed Of Light Readout]
  \label{def:physics:phyx-mini-0589:phyxminiproblems-problemphyxmini0589-vacuumspeedoflightreadout}
  \lean{PhyXMiniProblems.ProblemPhyXMini0589.vacuumSpeedOfLightReadout}
  Physlib's exact dimensionful vacuum light speed, read in named units
\end{definition}
\begin{proof}
  This is the declared model definition of vacuum Speed Of Light Readout.
\end{proof}
\begin{definition}[speed Fraction Of Light]
  \label{def:physics:phyx-mini-0589:phyxminiproblems-problemphyxmini0589-speedfractionoflight}
  \lean{PhyXMiniProblems.ProblemPhyXMini0589.speedFractionOfLight}
  \uses{def:physics:phyx-mini-0589:phyxminiproblems-problemphyxmini0589-speedquantity, def:physics:phyx-mini-0589:phyxminiproblems-problemphyxmini0589-speedreadout, def:physics:phyx-mini-0589:phyxminiproblems-problemphyxmini0589-vacuumspeedoflightreadout}
  A nonnegative speed expressed as a dimensionless fraction of vacuum c
\end{definition}
\begin{proof}
  This is the declared model definition of speed Fraction Of Light.
\end{proof}
\begin{definition}[velocity Fraction Of Light]
  \label{def:physics:phyx-mini-0589:phyxminiproblems-problemphyxmini0589-velocityfractionoflight}
  \lean{PhyXMiniProblems.ProblemPhyXMini0589.velocityFractionOfLight}
  \uses{def:physics:phyx-mini-0589:phyxminiproblems-problemphyxmini0589-signedvelocityquantity, def:physics:phyx-mini-0589:phyxminiproblems-problemphyxmini0589-signedvelocityreadout, def:physics:phyx-mini-0589:phyxminiproblems-problemphyxmini0589-vacuumspeedoflightreadout}
  A signed axial velocity expressed as a dimensionless fraction of c
\end{definition}
\begin{proof}
  This is the declared model definition of velocity Fraction Of Light.
\end{proof}
\begin{definition}[Inertial Frame Label]
  \label{def:physics:phyx-mini-0589:phyxminiproblems-problemphyxmini0589-inertialframelabel}
  \lean{PhyXMiniProblems.ProblemPhyXMini0589.InertialFrameLabel}
  The two inertial frames needed for the stated and pilot-observed speeds
\end{definition}
\begin{proof}
  This is the declared model definition of Inertial Frame Label.
\end{proof}
\begin{definition}[Cruiser Label]
  \label{def:physics:phyx-mini-0589:phyxminiproblems-problemphyxmini0589-cruiserlabel}
  \lean{PhyXMiniProblems.ProblemPhyXMini0589.CruiserLabel}
  The cruisers named in the prose and printed as A and B in the figure
\end{definition}
\begin{proof}
  This is the declared model definition of Cruiser Label.
\end{proof}
\begin{definition}[Figure Object]
  \label{def:physics:phyx-mini-0589:phyxminiproblems-problemphyxmini0589-figureobject}
  \lean{PhyXMiniProblems.ProblemPhyXMini0589.FigureObject}
  Physical objects explicitly distinguished by the primary image
\end{definition}
\begin{proof}
  This is the declared model definition of Figure Object.
\end{proof}
\begin{definition}[Cruiser Color]
  \label{def:physics:phyx-mini-0589:phyxminiproblems-problemphyxmini0589-cruisercolor}
  \lean{PhyXMiniProblems.ProblemPhyXMini0589.CruiserColor}
  The two cruiser colors visible in the bitmap
\end{definition}
\begin{proof}
  This is the declared model definition of Cruiser Color.
\end{proof}
\begin{definition}[Horizontal Direction]
  \label{def:physics:phyx-mini-0589:phyxminiproblems-problemphyxmini0589-horizontaldirection}
  \lean{PhyXMiniProblems.ProblemPhyXMini0589.HorizontalDirection}
  Direction along the left-to-right horizontal axis
\end{definition}
\begin{proof}
  This is the declared model definition of Horizontal Direction.
\end{proof}
\begin{definition}[Velocity Arrow Label]
  \label{def:physics:phyx-mini-0589:phyxminiproblems-problemphyxmini0589-velocityarrowlabel}
  \lean{PhyXMiniProblems.ProblemPhyXMini0589.VelocityArrowLabel}
  Mathematical labels printed over the two velocity arrows
\end{definition}
\begin{proof}
  This is the declared model definition of Velocity Arrow Label.
\end{proof}
\begin{definition}[Coordinate Axis Label]
  \label{def:physics:phyx-mini-0589:phyxminiproblems-problemphyxmini0589-coordinateaxislabel}
  \lean{PhyXMiniProblems.ProblemPhyXMini0589.CoordinateAxisLabel}
  The coordinate label printed at the right end of the horizontal axis
\end{definition}
\begin{proof}
  This is the declared model definition of Coordinate Axis Label.
\end{proof}
\begin{definition}[Cruiser Station Figure]
  \label{def:physics:phyx-mini-0589:phyxminiproblems-problemphyxmini0589-cruiserstationfigure}
  \lean{PhyXMiniProblems.ProblemPhyXMini0589.CruiserStationFigure}
  \uses{def:physics:phyx-mini-0589:phyxminiproblems-problemphyxmini0589-cruiserlabel, def:physics:phyx-mini-0589:phyxminiproblems-problemphyxmini0589-figureobject, def:physics:phyx-mini-0589:phyxminiproblems-problemphyxmini0589-cruisercolor, def:physics:phyx-mini-0589:phyxminiproblems-problemphyxmini0589-horizontaldirection, def:physics:phyx-mini-0589:phyxminiproblems-problemphyxmini0589-velocityarrowlabel, def:physics:phyx-mini-0589:phyxminiproblems-problemphyxmini0589-coordinateaxislabel}
  Typed qualitative evidence from phyx\_data/test\_image/589.png. Horizontal coordinates preserve visual ordering only and are not physical distances or a quantitative position scale
\end{definition}
\begin{proof}
  This is the declared model definition of Cruiser Station Figure.
\end{proof}
\begin{definition}[Equal Approach Cruiser Setup]
  \label{def:physics:phyx-mini-0589:phyxminiproblems-problemphyxmini0589-equalapproachcruisersetup}
  \lean{PhyXMiniProblems.ProblemPhyXMini0589.EqualApproachCruiserSetup}
  \uses{def:physics:phyx-mini-0589:phyxminiproblems-problemphyxmini0589-speedquantity, def:physics:phyx-mini-0589:phyxminiproblems-problemphyxmini0589-signedvelocityquantity, def:physics:phyx-mini-0589:phyxminiproblems-problemphyxmini0589-inertialframelabel, def:physics:phyx-mini-0589:phyxminiproblems-problemphyxmini0589-horizontaldirection, def:physics:phyx-mini-0589:phyxminiproblems-problemphyxmini0589-cruiserstationfigure}
  All physical observables in the scenario. In particular, cruiserBSpeedRelativeToStation is an independent dimensionful quantity: it is not defined from an answer choice or from the solved value 1/2
\end{definition}
\begin{proof}
  This is the declared model definition of Equal Approach Cruiser Setup.
\end{proof}
\begin{definition}[Matches Primary Figure]
  \label{def:physics:phyx-mini-0589:phyxminiproblems-problemphyxmini0589-matchesprimaryfigure}
  \lean{PhyXMiniProblems.ProblemPhyXMini0589.MatchesPrimaryFigure}
  \uses{def:physics:phyx-mini-0589:phyxminiproblems-problemphyxmini0589-equalapproachcruisersetup}
  The primary bitmap places the red cruiser A left of the blue cruiser B, and places both left of the station. Both labeled arrows point right. The axis and station markings are retained without assigning physical distances to the drawing
\end{definition}
\begin{proof}
  This is the declared model definition of Matches Primary Figure.
\end{proof}
\begin{definition}[Matches Problem Data]
  \label{def:physics:phyx-mini-0589:phyxminiproblems-problemphyxmini0589-matchesproblemdata}
  \lean{PhyXMiniProblems.ProblemPhyXMini0589.MatchesProblemData}
  \uses{def:physics:phyx-mini-0589:phyxminiproblems-problemphyxmini0589-speedfractionoflight, def:physics:phyx-mini-0589:phyxminiproblems-problemphyxmini0589-equalapproachcruisersetup}
  Frame assignments, physical directions, and the sole numerical input supplied by the prose. No station-frame speed is supplied for cruiser B
\end{definition}
\begin{proof}
  This is the declared model definition of Matches Problem Data.
\end{proof}
\begin{definition}[Has Physical Speed Parameters]
  \label{def:physics:phyx-mini-0589:phyxminiproblems-problemphyxmini0589-hasphysicalspeedparameters}
  \lean{PhyXMiniProblems.ProblemPhyXMini0589.HasPhysicalSpeedParameters}
  \uses{def:physics:phyx-mini-0589:phyxminiproblems-problemphyxmini0589-vacuumspeedoflightreadout, def:physics:phyx-mini-0589:phyxminiproblems-problemphyxmini0589-speedfractionoflight, def:physics:phyx-mini-0589:phyxminiproblems-problemphyxmini0589-equalapproachcruisersetup}
  Positivity and subluminality of the physical speeds. These domain conditions exclude the superluminal algebraic root but do not determine cruiser B's speed within the physical interval
\end{definition}
\begin{proof}
  This is the declared model definition of Has Physical Speed Parameters.
\end{proof}
\begin{definition}[Represents Approach Directions]
  \label{def:physics:phyx-mini-0589:phyxminiproblems-problemphyxmini0589-representsapproachdirections}
  \lean{PhyXMiniProblems.ProblemPhyXMini0589.RepresentsApproachDirections}
  \uses{def:physics:phyx-mini-0589:phyxminiproblems-problemphyxmini0589-velocityfractionoflight, def:physics:phyx-mini-0589:phyxminiproblems-problemphyxmini0589-equalapproachcruisersetup}
  The signs that give the word "approach" its directional meaning in B's frame: A, behind B, moves in the positive direction toward B, while the station, ahead of B, moves in the negative direction toward B
\end{definition}
\begin{proof}
  This is the declared model definition of Represents Approach Directions.
\end{proof}
\begin{definition}[Satisfies Collinear Einstein Velocity Transformations]
  \label{def:physics:phyx-mini-0589:phyxminiproblems-problemphyxmini0589-satisfiescollineareinsteinvelocitytransformations}
  \lean{PhyXMiniProblems.ProblemPhyXMini0589.SatisfiesCollinearEinsteinVelocityTransformations}
  \uses{def:physics:phyx-mini-0589:phyxminiproblems-problemphyxmini0589-speedfractionoflight, def:physics:phyx-mini-0589:phyxminiproblems-problemphyxmini0589-velocityfractionoflight, def:physics:phyx-mini-0589:phyxminiproblems-problemphyxmini0589-equalapproachcruisersetup}
  The collinear Einstein velocity transformation, written in dimensionless fractions of c. Since both cruisers move right in the station frame, beta\_A|B = (beta\_A - beta\_B) / (1 - beta\_A * beta\_B). The station has zero station-frame velocity, so its velocity in B's frame is -beta\_B. These are symbolic laws involving the still-unknown speed of B; neither equation contains its solved numerical value
\end{definition}
\begin{proof}
  This is the declared model definition of Satisfies Collinear Einstein Velocity Transformations.
\end{proof}
\begin{definition}[Approach Speeds Are Velocity Magnitudes]
  \label{def:physics:phyx-mini-0589:phyxminiproblems-problemphyxmini0589-approachspeedsarevelocitymagnitudes}
  \lean{PhyXMiniProblems.ProblemPhyXMini0589.ApproachSpeedsAreVelocityMagnitudes}
  \uses{def:physics:phyx-mini-0589:phyxminiproblems-problemphyxmini0589-speedfractionoflight, def:physics:phyx-mini-0589:phyxminiproblems-problemphyxmini0589-velocityfractionoflight, def:physics:phyx-mini-0589:phyxminiproblems-problemphyxmini0589-equalapproachcruisersetup}
  Operational interpretation of each pilot-reported nonnegative approach speed as the magnitude of the corresponding signed B-frame velocity
\end{definition}
\begin{proof}
  This is the declared model definition of Approach Speeds Are Velocity Magnitudes.
\end{proof}
\begin{definition}[Pilot Sees Equal Approach Speeds]
  \label{def:physics:phyx-mini-0589:phyxminiproblems-problemphyxmini0589-pilotseesequalapproachspeeds}
  \lean{PhyXMiniProblems.ProblemPhyXMini0589.PilotSeesEqualApproachSpeeds}
  \uses{def:physics:phyx-mini-0589:phyxminiproblems-problemphyxmini0589-equalapproachcruisersetup}
  The requirement posed by the question: the two independent physical approach speed observables measured by the pilot of B are equal. This premise does not give either speed a numerical value and does not assign an answer choice
\end{definition}
\begin{proof}
  This is the declared model definition of Pilot Sees Equal Approach Speeds.
\end{proof}
\begin{definition}[Answer Choice]
  \label{def:physics:phyx-mini-0589:phyxminiproblems-problemphyxmini0589-answerchoice}
  \lean{PhyXMiniProblems.ProblemPhyXMini0589.AnswerChoice}
  Labels of the four choices printed with the source question
\end{definition}
\begin{proof}
  This is the declared model definition of Answer Choice.
\end{proof}
\begin{definition}[displayed Cruiser BSpeed Fraction]
  \label{def:physics:phyx-mini-0589:phyxminiproblems-problemphyxmini0589-displayedcruiserbspeedfraction}
  \lean{PhyXMiniProblems.ProblemPhyXMini0589.displayedCruiserBSpeedFraction}
  \uses{def:physics:phyx-mini-0589:phyxminiproblems-problemphyxmini0589-answerchoice}
  Candidate station-frame speeds of cruiser B, as fractions of c
\end{definition}
\begin{proof}
  This is the declared model definition of displayed Cruiser BSpeed Fraction.
\end{proof}
\begin{definition}[recorded Dataset Answer]
  \label{def:physics:phyx-mini-0589:phyxminiproblems-problemphyxmini0589-recordeddatasetanswer}
  \lean{PhyXMiniProblems.ProblemPhyXMini0589.recordedDatasetAnswer}
  \uses{def:physics:phyx-mini-0589:phyxminiproblems-problemphyxmini0589-answerchoice}
  The answer label recorded by the dataset; it is not a theorem premise
\end{definition}
\begin{proof}
  This is the declared model definition of recorded Dataset Answer.
\end{proof}
\begin{definition}[Matches Answer Choice]
  \label{def:physics:phyx-mini-0589:phyxminiproblems-problemphyxmini0589-matchesanswerchoice}
  \lean{PhyXMiniProblems.ProblemPhyXMini0589.MatchesAnswerChoice}
  \uses{def:physics:phyx-mini-0589:phyxminiproblems-problemphyxmini0589-speedfractionoflight, def:physics:phyx-mini-0589:phyxminiproblems-problemphyxmini0589-equalapproachcruisersetup, def:physics:phyx-mini-0589:phyxminiproblems-problemphyxmini0589-answerchoice, def:physics:phyx-mini-0589:phyxminiproblems-problemphyxmini0589-displayedcruiserbspeedfraction}
  A choice exactly matches the modeled station-frame speed of cruiser B
\end{definition}
\begin{proof}
  This is the declared model definition of Matches Answer Choice.
\end{proof}
\begin{definition}[Is Unique Matching Answer Choice]
  \label{def:physics:phyx-mini-0589:phyxminiproblems-problemphyxmini0589-isuniquematchinganswerchoice}
  \lean{PhyXMiniProblems.ProblemPhyXMini0589.IsUniqueMatchingAnswerChoice}
  \uses{def:physics:phyx-mini-0589:phyxminiproblems-problemphyxmini0589-equalapproachcruisersetup, def:physics:phyx-mini-0589:phyxminiproblems-problemphyxmini0589-answerchoice, def:physics:phyx-mini-0589:phyxminiproblems-problemphyxmini0589-matchesanswerchoice}
  The selected choice is the unique displayed exact speed fraction
\end{definition}
\begin{proof}
  This is the declared model definition of Is Unique Matching Answer Choice.
\end{proof}
\begin{lemma}[required cruiser B speed fraction]
  \label{lem:physics:phyx-mini-0589:phyxminiproblems-problemphyxmini0589-required-cruiser-b-speed-fraction}
  \lean{PhyXMiniProblems.ProblemPhyXMini0589.required_cruiser_B_speed_fraction}
  \uses{def:physics:phyx-mini-0589:phyxminiproblems-problemphyxmini0589-speedfractionoflight, def:physics:phyx-mini-0589:phyxminiproblems-problemphyxmini0589-equalapproachcruisersetup, def:physics:phyx-mini-0589:phyxminiproblems-problemphyxmini0589-matchesproblemdata, def:physics:phyx-mini-0589:phyxminiproblems-problemphyxmini0589-hasphysicalspeedparameters, def:physics:phyx-mini-0589:phyxminiproblems-problemphyxmini0589-representsapproachdirections, def:physics:phyx-mini-0589:phyxminiproblems-problemphyxmini0589-satisfiescollineareinsteinvelocitytransformations, def:physics:phyx-mini-0589:phyxminiproblems-problemphyxmini0589-approachspeedsarevelocitymagnitudes, def:physics:phyx-mini-0589:phyxminiproblems-problemphyxmini0589-pilotseesequalapproachspeeds}
  Substituting beta\_A = 4/5 into the Einstein transformation and equating the two pilot-observed approach speeds gives the physical solution beta\_B = 1/2. The second algebraic root is superluminal and is excluded by HasPhysicalSpeedParameters
\end{lemma}
\begin{proof}
  The conclusion follows from the governing relations and figure readouts named in the statement of required cruiser B speed fraction.
\end{proof}
% --- Archon declaration topology end ---
% --- Archon physics formalization source end ---
